\section*{Chapter \ref{ch:pl:databases-and-sql} --- Databases and SQL}

\begin{solution}{exo:pl:databases-and-sql:1}
A positions table is a second copy of what the trades say, and two copies can disagree; a query over trades cannot. When the query is too slow,
a positions table can be kept as a cache, but it is rebuilt from the trades and never corrected by itself.
\end{solution}

\begin{solution}{exo:pl:databases-and-sql:2}
Because a corrected trade has several versions, one per system-time interval: the trade identifier repeats, and the start of \omterm{def:pl:databases-and-sql:system}{system time}
tells the versions apart.
\end{solution}

\begin{solution}{exo:pl:databases-and-sql:3}
500\,048 versions for 500\,008 trades: each of the 30 quantity corrections and 10 cancellations closes a version and adds one; the 8 late
bookings are new trades with one version each.
\end{solution}

\begin{solution}{exo:pl:databases-and-sql:4}
Of the six orders of two reads and two writes that keep each writer's read before its write, four put both reads before either write; then
both writers compute from the same starting value and the second write overwrites the first. The model's 636 of 1\,000 is close to four sixths.
\end{solution}

\begin{solution}{exo:pl:databases-and-sql:5}
Tuesday 18:00 in valid time selects every trade of Monday and Tuesday, two days of five or 40\% of the rows; reading them through the index
costs about as much as scanning the table.
\end{solution}

\begin{solution}{exo:pl:databases-and-sql:6}
Only the corrected report: the late trade is valid from Tuesday but its \omterm{def:pl:databases-and-sql:system}{system time} starts on Wednesday, so on Tuesday at 18:00 the database
did not hold it.
\end{solution}

\begin{solution}{exo:pl:databases-and-sql:7}
Join the trades table to itself on the trade identifier: one side restricted to the versions current at \omterm{def:pl:databases-and-sql:system}{system time} Tuesday 18:00, the other to
the versions current now, keeping pairs whose quantity or valid-time end differ. On the model's week it returns 40 trades: the 30 corrected
quantities and the 10 cancellations (query \texttt{CHANGED} in \texttt{pl\_tradedb.py}).
\end{solution}

\begin{solution}{exo:pl:databases-and-sql:8}
An audit table records the changes the application chose to log, in its own format; rebuilding a past state means replaying them backwards
without error, including changes made outside the application (a script, a manual fix). With \omterm{def:pl:databases-and-sql:system}{system time} the past states are the data
itself, queried like the present, and nothing can change them.
\end{solution}

\begin{solution}{pb:pl:databases-and-sql:1}
\begin{enumerate}
\item That each instrument, account and trade version is identified once, and that no trade refers to an account or instrument that does not
exist.
\item Facts about instruments and accounts (currency, names): they live once in their own tables.
\item So that every copy of the database applies the same changes once, in order, and can say which schema it has.
\item A valid-time interval and a system-time interval.
\item Corrected: the current version's \omterm{def:pl:databases-and-sql:system}{system time} is closed and the corrected version inserted. Cancelled: the same, with an empty valid-time
interval.
\item T1 reads an item, T2 updates it, and T1 writes a value computed from its earlier read and commits: T2's update is lost.
\item Because its read and its write are separate transactions: each is serializable, but the pair is not one transaction.
\item 636 read-then-write; none inside a transaction; none as one statement.
\item It takes the write lock at the start, so the second writer is refused until the first commits and then retries with the new value.
\item When the new value is a function of the old one that SQL can express, such as an increment.
\item A scan of the whole trades table for all three.
\item The (account, instrument) index takes the point query from 13.7~ms to 0.01~ms and leaves the others unchanged; the valid-time index
barely changes any.
\item DuckDB: report 9.8~ms against 150, analytical query 2.8~ms against 148, point query 1.0~ms against 13.7 without index and 0.01 with it.
\item Every insert and every correction must update it too, and it takes space.
\item A columnar copy (DuckDB here), refreshed from the transactional store.
\item \textbf{Named result.} Tuesday's 18:00 positions report is reproduced exactly from the bitemporal table (valid Tuesday 18:00, system
Tuesday 18:00) and differs from the corrected report (system Friday) on 48 of 10\,000 positions, one per correction: 30 quantities, 10
cancellations and 8 late bookings. The (account, instrument) index speeds the point query about a thousandfold (13.7~ms to 0.01~ms); the
valid-time index gains almost nothing on any query.
\item \omterm{def:pl:databases-and-sql:system}{System time} is stamped by the database when it records a fact and cannot be backdated; knowledge time is when the fact became available,
supplied with the data.
\item What the database believed before each correction, and so every report produced before it.
\item That it can answer as of any past \omterm{def:pl:databases-and-sql:system}{system time}.
\item What it believed, and when.
\end{enumerate}
\end{solution}

\begin{solution}{iq:pl:databases-and-sql:1}
A group of reads and writes applied as a unit: atomic (all or nothing), consistent (valid state to valid state), isolated (no partial work
seen by others) and durable (committed means it survives a crash).
\iqlookfor{All four, with what each protects.}
\end{solution}

\begin{solution}{iq:pl:databases-and-sql:2}
Both read the old value and write their own sum, and one write overwrites the other. Send increments as one statement, or do the read and the
write in one transaction that takes the write lock at the start, with a retry.
\iqlookfor{Read-modify-write in one transaction.}
\end{solution}

\begin{solution}{iq:pl:databases-and-sql:3}
Read the plan: a full scan of a large table for a query that selects few rows wants an index on the filtered columns; a query that reads most
rows does not. Measure before and after, and count the cost on writes.
\iqlookfor{Selectivity, plan, measurement.}
\end{solution}

\begin{solution}{iq:pl:databases-and-sql:4}
Keep \omterm{def:pl:databases-and-sql:system}{system time}: corrections close versions and insert new ones, never update; produce reports from queries as of a \omterm{def:pl:databases-and-sql:system}{system time}, and record
the \omterm{def:pl:databases-and-sql:system}{system time} each report used.
\iqlookfor{Bitemporal tables, never update in place.}
\end{solution}

\begin{solution}{iq:pl:databases-and-sql:5}
A transactional store with keys, migrations and bitemporal trades for booking; positions as queries or rebuilt caches; a columnar copy for
reporting and research, refreshed from it; indexes chosen from plans; every report stamped with its \omterm{def:pl:databases-and-sql:system}{system time}.
\iqlookfor{OLTP store of record, OLAP copy, time travel.}
\end{solution}
